% Shim de compatibilité pour les exercices MathAData
% Ce fichier est inséré tel quel dans le préambule du .tex généré
% (voir mathadataCompatTex dans src/lib/Latex.ts) : ce n'est pas un
% paquet chargé via \usepackage, donc pas de \NeedsTeXFormat ni de
% \ProvidesPackage ni de \RequirePackage ici, uniquement du \usepackage
% classique de préambule.

%%%%%%%%%%%%%%%%%%%%%%%%%%%%%%%%%%%%%%%%%%%%%%%%%%%%%%%%%%%%%%%%%%%%%%%%
% 1) Dépendances
%    xcolor et enumitem sont déjà chargés par le préambule générique de
%    Mathaléa (src/lib/latex/preambule.tex) : on ne les recharge pas ici.
%%%%%%%%%%%%%%%%%%%%%%%%%%%%%%%%%%%%%%%%%%%%%%%%%%%%%%%%%%%%%%%%%%%%%%%%
\usepackage{subcaption}         % fournit l'environnement subfigure (utilisé par les exercices MathAData)
\usepackage{circledsteps}       % fournit \Circled  (utilisé par MADenumerate variante A)
\usepackage{dashrule}           % fournit \hdashrule (utilisé par \ligne)
\usepackage[tikz]{bclogo}       % boîtes "exemple" / "bilan"

%%%%%%%%%%%%%%%%%%%%%%%%%%%%%%%%%%%%%%%%%%%%%%%%%%%%%%%%%%%%%%%%%%%%%%%%
% 2) Palette de charte
% Ces couleurs sont RÉFÉRENCÉES dans les figures des exercices mais
% jamais définies dans les fichiers -> sans elles, la compilation
% échoue (\color{BleuCeladon} indéfini). Aucun impact sur la mise en
% page. \providecolor évite tout conflit si une couleur existe déjà.
%%%%%%%%%%%%%%%%%%%%%%%%%%%%%%%%%%%%%%%%%%%%%%%%%%%%%%%%%%%%%%%%%%%%%%%%
\providecolor{CultureMath}{RGB}{230,29,73}
\providecolor{BleuCeladon}{RGB}{20,129,173}
\providecolor{VertMenthe}{RGB}{62,178,136}
\providecolor{JauneMoutarde}{RGB}{250,181,29}
\providecolor{OrangeCarotte}{RGB}{241,132,18}
\providecolor{RougeFramboise}{RGB}{230,29,73}
\providecolor{RosePerle}{RGB}{244,167,198}
\providecolor{GrisDeSoutien}{RGB}{245,250,248}
\providecolor{floralwhite}{RGB}{255,250,240}

%%%%%%%%%%%%%%%%%%%%%%%%%%%%%%%%%%%%%%%%%%%%%%%%%%%%%%%%%%%%%%%%%%%%%%%%
% 3) \ligne  — trait pointillé à compléter par l'élève
%    Usage : \ligne  ou  \ligne[4cm]
%%%%%%%%%%%%%%%%%%%%%%%%%%%%%%%%%%%%%%%%%%%%%%%%%%%%%%%%%%%%%%%%%%%%%%%%
\providecommand{\ligne}[1][2.8cm]{\hdashrule{#1}{0.5pt}{0.5pt 3pt}\ }

%%%%%%%%%%%%%%%%%%%%%%%%%%%%%%%%%%%%%%%%%%%%%%%%%%%%%%%%%%%%%%%%%%%%%%%%
% 4) MADenumerate — liste numérotée à pastilles
%
%    Repère visuel volontairement conservé : une pastille numérotée
%    aide l'élève à se situer dans les questions. C'est la FORME
%    (pastille) qui porte l'aide pédagogique, pas la couleur.
%
%    La couleur est donc paramétrable : Mathaléa peut l'aligner sur sa
%    propre couleur d'accent en redéfinissant \MADbulletcolor, sans
%    rien changer d'autre :
%        \renewcommand{\MADbulletcolor}{black}   % ou une couleur maison
%%%%%%%%%%%%%%%%%%%%%%%%%%%%%%%%%%%%%%%%%%%%%%%%%%%%%%%%%%%%%%%%%%%%%%%%
\providecommand{\MADbulletcolor}{CultureMath}

% ----- VARIANTE A (recommandée) : pastille numérotée -----------------
% Le \setlist[...,2] reste LOCAL au groupe de l'environnement : il ne
% modifie pas le enumerate global de Mathaléa.
\newenvironment{MADenumerate}[1][1]
 {\setlist[enumerate,2]{leftmargin=0.3em, itemindent=1.2em, labelsep=0.3em}%
  \begin{enumerate}[%
      label={\Circled[fill color=\MADbulletcolor, outer color=white, inner color=white]{\arabic*}},%
      font=\large\bfseries, start=#1,%
      leftmargin=0.2em, itemindent=1.3em, labelsep=0.5em]}
 {\end{enumerate}}

% ----- VARIANTE B : alias neutre vers enumerate ----------------------
% Adopte le style de liste natif de Mathaléa, aucune dépendance
% (ni circledsteps, ni couleur). Décommenter B et commenter A pour
% basculer.
%
% \newenvironment{MADenumerate}[1][1]
%  {\begin{enumerate}[start=#1]}
%  {\end{enumerate}}

%%%%%%%%%%%%%%%%%%%%%%%%%%%%%%%%%%%%%%%%%%%%%%%%%%%%%%%%%%%%%%%%%%%%%%%%
% 5) Boîtes sémantiques "exemple" et "bilan"  (2 usages chacune)
%    Version fidèle (bclogo). Peu d'enjeu : Mathaléa peut les
%    redéfinir en boîtes maison si elles jurent avec leur charte.
%%%%%%%%%%%%%%%%%%%%%%%%%%%%%%%%%%%%%%%%%%%%%%%%%%%%%%%%%%%%%%%%%%%%%%%%
\newenvironment{exemple}[1][Exemple]
 {\begin{bclogo}[epBord=0.5, arrondi=0.1, logo=\bcloupe,
                 couleur=JauneMoutarde!20, ombre=true, epOmbre=0.2,
                 noborder=true]{#1 }}
 {\end{bclogo}\vspace{-10pt}}

\newenvironment{bilan}[1][Bilan]
 {\begin{bclogo}[epBord=0.5, arrondi=0.1, logo=\bccrayon,
                 couleur=VertMenthe!20, ombre=true, epOmbre=0]{#1}}
 {\end{bclogo}\vspace{-10pt}}
